% Introducción

\nnchapter{INTRODUCCIÓN} \label{chap:introduccion}

La diabetes mellitus (DM) es una enfermedad metabólica crónica que se caracteriza por
niveles elevados de glucosa en la sangre. Esto ocurre cuando el cuerpo no produce
suficiente insulina, una hormona producida por el páncreas que regula la cantidad de
glucosa en la sangre, o cuando las células no responden adecuadamente a la insulina.
Alrededor de 463 millones de adultos entre 20 y 79 años tienen diabetes mellitus (DM) en
la actualidad, lo que representa el 9.3\% de la población mundial en este grupo de edad
\myfootcite{antonetti2021current}.

La retinopatía diabética (RD) es una complicación
microvascular específica de la diabetes mellitus (DM) que afecta a 1 de cada 3 personas
con esta enfermedad, dañando los vasos sanguíneos de la retina y convirtiéndose en una de
las principales causas de pérdida de visión en la población adulta activa. Los pacientes
con niveles graves de RD experimentan una calidad de vida significativamente reducida y
niveles más bajos de bienestar físico, emocional y social. Frente a este panorama, el 
presente trabajo propone el desarrollo de un algoritmo de
clasificación basado en aprendizaje profundo, específicamente mediante redes neuronales
convolucionales (CNN), para la detección de retinopatía diabética en imágenes del fondo
del ojo. 

Este proyecto utiliza estas redes en la construcción de clasificadores binarios
(presencia o ausencia de la enfermedad) y multiclase (cinco niveles de severidad), dado que
este enfoque combinado ha mostrado obtener resultados prometedores tanto para tamizaje como
para gradación clínica. El modelo binario permite una rápida detección de la enfermedad,
optimizando el tiempo de diagnóstico y brindando una solución viable para su implementación en
escenarios de alta demanda, como en áreas con recursos limitados. Por su parte, los
clasificadores multiclase aportan valor adicional en términos de tratamiento personalizado y
seguimiento de los pacientes al distinguir entre los diferentes grados de severidad de la retinopatía
diabética.

Una de las principales contribuciones de este proyecto es la experimentación con diferentes
configuraciones de redes neuronales, explorando cómo las variaciones en el preprocesamiento
y las arquitecturas afectan el desempeño del modelo. Se implementaron y compararon cuatro
arquitecturas preentrenadas: MobileNetV2, ResNet50V2, InceptionV3 y VGG19, utilizando
aprendizaje por transferencia con 35.121 imágenes del dataset de Kaggle para entrenamiento
y datasets externos (MESSIDOR-2, IDRiD, FOSCAL) para validación. Un aspecto distintivo es
el análisis de imágenes a través de los tres canales de color RGB (rojo, verde y azul) de
forma independiente, un enfoque que no ha sido ampliamente explorado en investigaciones
previas. La idea es determinar cuál de estos canales, por separado, proporciona la mejor
representación de las características relevantes para la detección de la RD. Este análisis
permitió identificar el canal verde como el más efectivo para la clasificación multiclase
debido a su mejor contraste para lesiones hemorrágicas.

Adicionalmente, se implementaron distintas técnicas de regularización, como Dropout,
regularización L2, capas GaussianNoise, bloques Multi-Head Attention y Data Augmentation,
con el fin de reducir el sobreajuste y mejorar la generalización del modelo. Estas técnicas
se evaluaron de manera exhaustiva para identificar cuáles son más efectivas en el contexto
específico de clasificación de imágenes de retinopatía diabética. La combinación de estas
estrategias no solo mejoró el rendimiento del modelo, sino que también contribuyó a
desarrollar un sistema robusto y adaptable a diferentes tipos de datos, mejorando su
aplicabilidad en diversas poblaciones. Este trabajo logró una sensibilidad de 96.5\%
en clasificación binaria (con intervalos de confianza bootstrap del 95\%: AUC-ROC 0.9636–0.9936)
y 65\% de sensibilidad en clasificación multiclase, demostrando la viabilidad de estas
tecnologías para tamizaje automatizado.

Esta investigación busca impulsar el uso de tecnologías recientes en salud ocular. Su
aplicación permitirá abordar la falta de especialistas en oftalmología, promover la
detección oportuna de enfermedades y ampliar el alcance del diagnóstico, optimizando
así la eficiencia y accesibilidad de los servicios de salud visual. Además, esta
metodología podría sentar las bases para implementar soluciones similares en otros países
en vías de desarrollo, generando un impacto significativo en la calidad de vida de los
pacientes.

El presente documento se estructura de la siguiente manera: en el Capítulo 1 se plantea
el problema que motiva este trabajo; el Capítulo 2 detalla los objetivos; los Capítulos 3
y 4 presentan el marco teórico y los antecedentes; el Capítulo 5 describe la metodología
implementada; el Capítulo 6 expone los resultados obtenidos; finalmente, los Capítulos 7,
8 y 9 presentan la discusión, conclusiones y recomendaciones respectivamente.

